\begin{trapbox}
\begin{description}[style=nextline, leftmargin=2.8em, labelsep=0.5em, itemsep=0.35em]
    \item[\textbf{\textcolor{CTrap}{易错点一（忽视正数条件）}}] 直接对负数或可能为负的变量使用 $a + b \ge 2\sqrt{ab}$。
    \item[\textbf{\textcolor{CTrap}{正确理解（正数前提）：}}] 公式 $a + b \ge 2\sqrt{ab}$ 和 $ab \le ((a+b)/2)^2$ 仅当 $a, b \ge 0$ 时成立。
    \item[\textbf{\textcolor{CTrap}{错因分析（根式无意义）：}}] 若 $a, b$ 为负，$\sqrt{ab}$ 可能无实数意义；即使有意义（如两负相乘为正），算术-几何平均不等式也不一定成立。
\end{description}

\medskip
\begin{description}[style=nextline, leftmargin=2.8em, labelsep=0.5em, itemsep=0.35em]
    \item[\textbf{\textcolor{CTrap}{易错点二（忽略取等条件）}}] 求出最值后，不验证等号是否能成立，导致最值不可达。
    \item[\textbf{\textcolor{CTrap}{正确理解（同时性检验）：}}] 基本不等式等号成立当且仅当 $a=b$。在复杂表达式中，需检查此条件是否能在约束下同时满足。
    \item[\textbf{\textcolor{CTrap}{错因分析（过程与结果脱节）：}}] 推导过程中可能多次使用不等式，若等号成立条件不一致，则最终等号可能无法同时取到，所求“最值”仅为下界而非实际最值。
\end{description}

\medskip
\begin{description}[style=nextline, leftmargin=2.8em, labelsep=0.5em, itemsep=0.35em]
    \item[\textbf{\textcolor{CTrap}{易错点三（混淆公式范围）}}] 认为 $a^2 + b^2 \ge 2ab$ 也要求 $a, b$ 为正数。
    \item[\textbf{\textcolor{CTrap}{正确理解（实数范围成立）：}}] $a^2 + b^2 \ge 2ab$ 由 $(a-b)^2 \ge 0$ 推出，对任意实数 $a, b$ 都成立。
    \item[\textbf{\textcolor{CTrap}{错因分析（记忆混淆）：}}] 将不同形式基本不等式的适用条件记混，未区分“算术-几何平均”与“平方和与积”两种形式的前提差异。
\end{description}
\end{trapbox}
